% =============================================
% CHƯƠNG 2 — KHẢO SÁT VÀ PHÂN TÍCH YÊU CẦU (9 trang)
% Status: OUTLINE — chờ draft từ agent
% =============================================

\chapter{Khảo sát và phân tích yêu cầu}
\label{ch:khao-sat}

% [INTRO CHƯƠNG — 1 đoạn]
% Chương này khảo sát các hệ thống điểm danh hiện có (trong và ngoài nước),
% phân tích ưu nhược điểm của từng cách tiếp cận, từ đó rút ra các yêu cầu
% chức năng và phi chức năng cho hệ thống cần xây dựng.

% ─────────────────────────────────────────────
\section{Khảo sát các hệ thống điểm danh hiện có}
\label{sec:khao-sat-he-thong}
% Ước lượng: 3 trang
% VIẾT MỚI nhiều, dùng IEEE refs
% Nội dung:
%   §2.1.1 Điểm danh truyền thống (gọi tên, ký giấy) — 0.5tr
%   §2.1.2 RFID / Thẻ từ — 0.5tr (cite paper RFID attendance system)
%   §2.1.3 NFC — 0.5tr
%   §2.1.4 BLE Beacon — 0.5tr (cite Crowley et al, Apple iBeacon)
%   §2.1.5 QR Code tĩnh — 0.5tr (vấn đề screenshot, share QR)
%   §2.1.6 Nhận diện khuôn mặt độc lập — 0.5tr (cite FaceNet)
%   * Mỗi mục: cơ chế hoạt động, ưu điểm, nhược điểm, ví dụ thực tế
%   * Tham chiếu: tìm 4-6 paper IEEE qua agent citation

\section{So sánh và bài học rút ra}
\label{sec:so-sanh}
% Ước lượng: 1.5 trang
% Nội dung: bảng so sánh đa chiều
%   - Tiêu chí: chi phí phần cứng, chi phí triển khai, độ chính xác, chống gian lận,
%     trải nghiệm SV, khả năng mở rộng, có cần app riêng không
%   - Bảng 2.1: ma trận 6 phương án × 7 tiêu chí
%   - Bài học: cần kết hợp nhiều layer, không phương án đơn lẻ nào đủ tốt
% Nguồn: báo cáo 13-18/04 mục "Nghiên cứu và tham khảo"

% ─────────────────────────────────────────────
\section{Phân tích đối tượng người dùng}
\label{sec:doi-tuong}
% Ước lượng: 1 trang
% Nội dung:
%   §2.3.1 Sinh viên — đặc điểm, kỳ vọng, ràng buộc (tốc độ <10s, không tốn data)
%   §2.3.2 Giảng viên — đặc điểm, kỳ vọng (giám sát realtime, xuất báo cáo)
%   §2.3.3 Phòng Đào tạo / Admin — đặc điểm (quản lý hàng loạt, import Excel, thống kê)

% ─────────────────────────────────────────────
\section{Yêu cầu chức năng}
\label{sec:yeu-cau-cn}
% Ước lượng: 1.5 trang
% Nguồn: src/types/index.ts + 16 services + Notion page 3+4+5
% Nội dung:
%   Bảng 2.2 — Yêu cầu chức năng theo 3 vai trò
%   Mỗi vai trò liệt kê 5-8 chức năng chính, đánh số FR-X.Y
%   - SV: FR-S.1 đăng nhập Zalo, FR-S.2 đăng ký khuôn mặt, FR-S.3 quét QR, ...
%   - GV: FR-T.1 tạo phiên, FR-T.2 giám sát realtime, FR-T.3 duyệt thủ công, ...
%   - Admin: FR-A.1 quản lý SV, FR-A.2 quản lý lớp, FR-A.3 duyệt đơn xin nghỉ, ...

% ─────────────────────────────────────────────
\section{Yêu cầu phi chức năng}
\label{sec:yeu-cau-pcn}
% Ước lượng: 1.5 trang
% Nguồn: SECURITY-REPORT.md
% Nội dung:
%   §2.5.1 Bảo mật — HMAC, OAuth, Firestore rules, encryption at rest
%   §2.5.2 Hiệu năng — cold start <3s, quét QR <1s, render screen <500ms
%   §2.5.3 Khả năng mở rộng — Firestore tự scale, target 5000 SV
%   §2.5.4 Tính sẵn sàng — uptime 99% (giới hạn bởi Firebase Spark)
%   §2.5.5 Trải nghiệm — UI tối ưu điện thoại Android+iOS, dark/light theme
%   §2.5.6 Tuân thủ — Zalo Mini App Censorship Policy, GDPR-like privacy
%   Bảng 2.3 — Tổng hợp NFR với metrics đo lường

% ─────────────────────────────────────────────
\section{Use case và Actor}
\label{sec:use-case}
% Ước lượng: 0.5 trang
% Nguồn: Notion page 3, 4, 5 (sequence + state diagrams)
% Nội dung:
%   - Hình 2.1: Sơ đồ use case tổng thể (mermaid → PNG)
%   - 3 Actors: SV, GV, Admin
%   - 15-20 use case chính (đặt tên UC-NN)
%   - Liệt kê use case theo bảng (không mô tả chi tiết — để chương 4)

% ─────────────────────────────────────────────
% \section*{Kết chương}
% Tóm tắt: đã khảo sát 6 phương án, rút ra cần kết hợp đa layer; đã định nghĩa được
% yêu cầu FR và NFR. Chương tiếp theo sẽ phân tích công nghệ cụ thể.
